\newgeometry{margin=.7cm}
\pagestyle{empty}
\raggedcolumns
\footnotesize
\newcommand{\cheatSheetColor}{RoyalBlue}
\newtcolorbox{cheatSheetBox}[1][]{
	enhanced jigsaw,
%	breakable,
	colback=\cheatSheetColor!10,
	colframe=\cheatSheetColor,
	colbacktitle=\cheatSheetColor!30,
	fonttitle=\bfseries,
	parbox=false,
	attach boxed title to top text left={yshift=-3mm,yshifttext=-3mm},
	coltitle=black,
	#1}

\begin{landscape}
\footnotesize
\begin{multicols*}{4}
\textcolor{\cheatSheetColor}{\begin{center}
\section*{Cheatsheet}
\end{center}\hrule}

\begin{cheatSheetBox}[title=Variablen erstellen]
\begin{center}
$ \tikzmarknode{b}{\text{\texttt{a}}} =\tikzmarknode{a}{\color{red}\underbrace{\color{black}\text{\texttt{10+3}}}}$
	
	\tikz[remember picture, overlay]{\draw[-latex,red] ([yshift=0.1em]a.south) to[bend left] ([yshift=-0.3em]b.south);}
\end{center}
\end{cheatSheetBox}

\begin{cheatSheetBox}[title=Schleifen mit \lstinline|for|]
\begin{lstlisting}
for _ in range(5):
	# mach etwas 5 mal!
# hier ist die Schleife fertig
\end{lstlisting}
\end{cheatSheetBox}

\begin{cheatSheetBox}[title=Modulo (\lstinline|%|)]
Modulo berechnet den Rest der Ganzzahldivision.
\begin{lstlisting}
print(9 % 5) # 4
\end{lstlisting}
\end{cheatSheetBox}


\begin{cheatSheetBox}[title=Wert bei Ausführung eingeben]
\begin{lstlisting}
name = input("name")
print("Hallo" + name)
\end{lstlisting}
\end{cheatSheetBox}

%\begin{cheatSheetBox}[title=Arbeiten mit Text]
%\end{cheatSheetBox}
\end{multicols*}

%\begin{table}[H]
% \centering
%\begin{tabular}{l|p{0.3\linewidth}|p{0.4\linewidth}}
%\textbf{Python} & \textbf{Beschreibung} & \textbf{Beispiel} \\\midrule\midrule
%\lstinline|len(s)| & Gibt die Länge des Textes \lstinline|s| zurück. & 
%
%\begin{minipage}{\linewidth}
%\begin{lstlisting}[numbers=none]
%n = len("hallo") # 5
%\end{lstlisting}
%\end{minipage} \\\midrule\midrule
%\lstinline|s[index]| & Greift auf das Zeichen an der Index-Position \lstinline|index| zu. Bei negativen Indizes zählt man von rechts. & \begin{minipage}{\linewidth}
%\begin{lstlisting}[numbers=none]
%s = "hallo"
%s = "hallo"
%c0 = s[0]  # "h"
%c1 = s[1]  # "a"
%c2 = s[-1] # "o"
%\end{lstlisting}
%\end{minipage}
% \\\midrule
%\lstinline|s[start:end]| & Extrahiert den Teiltext vom Index \lstinline|start| bis und ohne \lstinline|end| & \begin{minipage}{\linewidth}
%\begin{lstlisting}[numbers=none]
%s1 = "hallo"
%s2 = s1[1:4] # "all"
%\end{lstlisting}
%\end{minipage}\\\midrule\midrule
%\lstinline|ord(c)| & Gibt die Position des Zeichens \lstinline|c| in der \href{https://en.wikipedia.org/wiki/List_of_Unicode_characters}{Unicode-Tabelle} zurück & \begin{minipage}{\linewidth}
%\begin{lstlisting}[numbers=none]
%n1 = ord('A') # 65
%n2 = ord('B') # 66
%n3 = ord('Z') # 90
%\end{lstlisting}
%\end{minipage}\\\midrule
%\lstinline|chr(n)| & Gibt die das Zeichen an der \lstinline|n|-ten Position er \href{https://en.wikipedia.org/wiki/List_of_Unicode_characters}{Unicode-Tabelle} zurück & \begin{minipage}{\linewidth}
%\begin{lstlisting}[numbers=none]
%c1 = chr(65) # "A"
%c2 = chr(66) # "B"
%c3 = chr(67) # "C"
%\end{lstlisting}\end{minipage}\\\midrule\midrule
%\lstinline|s1.count(s2)| & Zählt, wie oft der Text \lstinline|s2| im Text \lstinline|s1| vorkommt & \begin{minipage}{\linewidth}
%\begin{lstlisting}[numbers=none]
%s1 = "Das Haus ist gross und das Dach ist grün."
%s2 = "ist"
%n = s1.count(s2) # 2
%\end{lstlisting}
%\end{minipage}\\\midrule
%\lstinline|s1.find(s2)| & Gibt den niedrigsten Index des Teiltexts \lstinline|s2| zurück, falls dieser vorhanden ist, ansonsten \lstinline|-1| & \begin{minipage}{\linewidth}
%\begin{lstlisting}[numbers=none]
%s = "Das Haus ist gross und das Dach ist grün."
%n = s.find("ist") # 9
%\end{lstlisting}
%\end{minipage}\\\midrule
%\lstinline|s1.replace(old, new)| & Ersetzt alle Vorkommen des Teil-Texts \lstinline|old| durch \lstinline|new|. & \begin{minipage}{\linewidth}
%\begin{lstlisting}[numbers=none]
%s1 = "hallo"
%s2 = s1.replace("a", "e") # "hello"
%\end{lstlisting}
%\end{minipage}
%\end{tabular}
%\caption{Übersicht von Zeichenketten-Befehlen}
%\label{tab:string-op}
%\end{table}


\end{landscape}

